\newpage
\section{PRESENTACIÓN}
% \section{\textcolor{mycolor}{PRESENTACIÓN}}
El propósito del presente informe es presentar al Comité de Operación Económica del Sistema Interconectado Nacional (COES-SINAC) el Estudio Hidrológico del Sistema Hídrico de ENGIE Energía Perú para el periodo 1965-2025, cumpliendo con lo estipulado en la Resolución OSINERGMIN No 196-2016-OS/CD del 02 de agosto del 2016, procedimiento técnico No 41. El Sistema Hídrico en estudio abarca la cuenca del río Paucartambo que es aprovechada con fines hidroeléctricos en la central hidroeléctrica de Yuncán - ENGIE Energía Perú S.A.\\

En la sección \ref{section:Methodology} se presenta la metodología seguida para el cálculo de caudales naturales, estimación de la precipitación satelital y su validación con las estaciones existentes en campo, identificación del número de estaciones virtuales (del satélite) de precipitación para la correcta caracterización por subcuencas y los análisis de las precipitaciones y de las series de caudales naturalizados. La sección 5 muestra la caracterización hidrológica de la cuenca del río Paucartambo hasta el embalse Yuncán (Mapa 5-1). La sección 6 presenta el modelo hidrológico distribuido \textbf{VIC (Variable Infiltration Capacity)}, sus insumos de entrada y los resultados de las simulaciones para obtener los caudales promedio mensuales desde 1998 al 2025 en los puntos QN-908 - Uchuhuerta, QN-909 – Huallamayo, QN-910 – Victoria, QN-901 + QN-902 + QN-911, y QN-903 + QN-904 + QN-905. Las salidas del modelo fueron corregidas mediante redes neuronales para ajustar los caudales simulados a los datos observados disponibles. La sección 7 muestra las discusiones y conclusiones.\\

El diagrama topológico se presenta en Mapa 4-1 junto con el mapa de la cuenca y la ubicación de las estaciones de precipitación existentes en la cuenca.\\

Se hace hincapié en los caudales mensuales naturalizados (Gráfico 6-1) y los mensuales promedios multianuales (Gráfico 6-2) en QN-908 – Uchuhuerta, los respectivos (Gráfico 6-3 y 6-4) en QN-909 – Huallamayo y para (Gráfico 6-5 y 6-6) QN-910 – Victoria I.

% Ya que la actual metodología (hidrología física) y los desarrollos tecnológicos (tecnología satelital - https://pmm.nasa.gov/GPM) permiten obtener caudales naturales directamente. No se han calculado los cambios en los embalses ni otra información adicional, ya que no se requiere. Con las metodologías antiguas, la estimación de caudales se derivaba en base a la información de los cambios en los embalses. En la actualidad, esto ya no es necesario.

\section{ANTECEDENTES}

El \textbf{COES-SINAC}, por ley, debe al \textbf{OSINERGMIN} los estudios económicos, información técnica, entre otra información necesaria para que ésta determine los precios en barra de potencia y energía del sistema interconectado. Para el cálculo de costos marginales del sistema interconectado, el \textbf{OSINERGMIN} necesita la información hidrológica de las cuencas que aprovechan las empresas generadoras.
